\section{Sistema normalizado de ecuaciones diferenciales ordinarias}
Al hacer las siguientes transformaciones sobre las ecuaciones \eqref{eq:partODE}  y \eqref{eq:energyODE} Haciendo siguiente cambio de variables para normalizar el problema:

\begin{equation}
	\rho = \frac{r}{R}, \quad \hat{n} = \frac{n - n_{min}}{\Delta_n}, \quad  \hat{T} = \frac{T - T_{min}}{\Delta_T},
\end{equation}

con $\Delta_x = x_{max} - x_{min}$. Se procedió a adimensionar las ODEs y remover posibles denominadores que pudieran presentar valores singulares durante la solución con la PINN, de tal forma que se llega a una ODEs normalizadas numéricamente estables. Para la ODE del flujo de partículas:

\begin{eqnarray}\label{eq:partODEnorm}
	&&\frac{n}{\Delta_n}\frac{T}{\Delta_T}\left[\rho\frac{d^2\hat{n}}{d\rho^2} + \frac{d\hat{n}}{d\rho}\right] + \rho\frac{T}{\Delta_T}\left(\frac{d\hat{n}}{d\rho}\right)^2- \frac{\rho}{2}\frac{n}{\Delta_n}\frac{d\hat{n}}{d\rho}\frac{d\hat{T}}{d\rho}+ \frac{R^2\Delta_T^{1/2}}{h\Delta_n^2}\left(\frac{T}{\Delta_T}\right)^{3/2}\rho N(\hat{n}, \hat{T}) = 0
\end{eqnarray}

donde se usa $x / \Delta_x =\hat{x} + x_{min}/\Delta_x$. De forma similar se llega a la expresión para la ODE del flujo de energía:

\begin{eqnarray}\label{eq:energyODEnorm}
	&& \frac{47}{10}\left(\frac{n}{\Delta_n}\right)^2\frac{T}{\Delta_T}\left(\rho\frac{d^2\hat{T}}{d\rho^2} + \frac{d\hat{T}}{d\rho}\right) +  \frac{52\rho}{5}\frac{n}{\Delta_n}\frac{T}{\Delta_T}\frac{d\hat{n}}{d\rho}\frac{d\hat{T}}{d\rho} -\frac{47\rho}{20}\left(\frac{n}{\Delta_n}\frac{d\hat{n}}{d\rho}\right)^2 - \frac{3w_b}{\Delta_T}\frac{T}{\Delta_T}\left(\frac{n}{\Delta_n}\right)^2 \nonumber\\ 
	&&- \frac{R^2\Delta_T^{1/2}}{h\Delta_n^2}\left(\frac{T}{\Delta_T}\right)^{3/2}\frac{\rho}{\Delta_T} \biggl\{TN(\hat{n},\hat{T})+S(\hat{n},\hat{T}) -\frac{P_{ext}}{\sqrt{2\pi}\sigma}\exp{\left[-\frac{1}{2}\left(\frac{R}{\sigma}\right)^2(\rho-\hat{\mu})^2\right]}\biggl\} = 0
\end{eqnarray}

con $w_b \equiv e^2R^2B^2 / m_i$ y $\hat{\mu} = \langle \rho \rangle$. Los parámetros libres del modelo son $\{\hat{\mu}, \sigma\}$ para el modelado del perfil Gaussiano de deposición de potencia. Mientras que los parámetros $\{n_{min}, T_{min}, n_{max}, T_{max}, P_{ext}, B, R\}$ se pueden ajustar con los datos típicos del dispositivo que se está estudiando.
